\subsection{共享缓存下的组合搜索与同方案对照}
问题三采用与问题二相同的有界结构候选族，并增加由实际 FIFO 时间线引导的缓存邻域。先在 L2 规则下重新评价声明的场景 B 种子和已有 L2 锚点，之后围绕重复读取的消费者、合法迁核与顺序插入生成候选。缓存候选按当前最优方案指纹组织有限菜单，不能直接更改 FIFO 淘汰规则或手工指定命中事件。

自适应阶段前 $16$ 次提议且处于预留探索时间的前 $40\%$ 时，交替尝试缓存邻域和普通联合邻域；随后先补充未探索的类别，再依据结构先验、实测收益和耗时分配机会。缓存类初始权重为 $2$，拓扑前沿为 $1$，存在独立分量时开放装箱类。总时限、最终复核预留、提议上限及生成准备复用与问题二一致。命中率只解释数据服务路径，最终仍按 $(T^{L2},D)$ 选择，不能以命中更多代替周期更短。

本轮初解是显式保存的历史输入：二至四核与五核的场景 B 种子来自已登记的不同证据批次，均在 L2 下重新评价；不是将本轮新问题二的结果自动传给问题三。因此，跨问最终成绩的比较只表示各自选定方案的综合差别，不代表单独的硬件收益。

为隔离硬件作用，对最终选定的同一问题三方案 $p_{L,i,k}$ 额外执行无 L2 官方重放，定义
\begin{equation}
R^{\mathrm{same}}_{i,k}=\frac{T^B_i(p_{L,i,k})}{T^{L2}_i(p_{L,i,k})},\qquad
\overline R^{\mathrm{same}}_k=\frac1{100}\sum_{i=1}^{100}R^{\mathrm{same}}_{i,k}.
\end{equation}
同时报告 $T^B_i(p_{B,i,k})/T^{L2}_i(p_{L,i,k})$ 的逐图均值作为跨方案综合比较。前者保持划分、核心与顺序不变，后者同时改变方案；两者不得互称。附录为最终问题三方案列出无 L2 与只读 L2 两种配置的周期、额外搬运量和命中率。相同方案两次评价的逻辑额外搬运应相等，服务池变化只影响物理传输及完成时刻。

题目前两问的单核加速比固定为 $1$；本问若展示相对无 L2 单核基准的 L2 曲线，单核点可以大于 $1$。该变化来自同一整图方案启用缓存，不能冒称单核并行收益。所有平均值均先按图计算比值，再对图取算术平均。
